\scanpage{6}
\mspage{2}
suivant :
\[
(4)\quad\left\{\begin{aligned}
&\text{a) Une courbe projective et lisse }Y(=\widehat{\widetilde X})\deleted{\unsure{sur }k}\\
&\text{b) Deux parties }\inserted{\text{finies}}\text{ disjointes }\widetilde A\text{ et }I\text{ de }Y(k)\\
&\text{(ou encore, deux sous-schémas finis étales}\\
&\text{disjoints }\widetilde A_k\text{ et }I_k\text{ de }Y\text{)}\\
&\text{c) Une involution sans point fixe }\sigma_{\widetilde A}\text{ de}\\
&\widetilde A\text{ (qui est donc de cardinal pair).}
\end{aligned}\right.
\]
Inversement, partant de données a) b) c), on construit une courbe standart $X$ en passant au quotient dans $Y\setminus I_k\supset\widetilde A_k$ par l'involution $\sigma$ -- i.e. $X$ est universel \inserted{(dans (Sch))} pour la donnée $p$ :
\diagram{\begin{tikzcd}
Y\setminus I_k\arrow[r,"p"] & X\\
\bigl(\,\widetilde A_k\arrow[u,"i"]\arrow[ur,dashed] & \bigr)
\end{tikzcd}}
, soumise à $(pi)\sigma=pi$.

Ainsi la catégorie des courbes standart sur $k$ \inserted{(pour les iso)} apparaît comme équivalente à celle des systèmes a) b) c) ci-dessus. (pour les iso)\ldots

\uline{NB} On récupère $\widehat X$ comme quotient de $Y$ par $\sigma$.
\mspage{5}
